% Teaching manuscript. Edit this TeX source; Markdown is compatibility navigation only.

\StudyChapter{Foundations-NetworkBasics}{Network Basics｜把单位、字节和时间轴读明白}
\StudyTarget{Foundations-NetworkBasics__network-basics}


\textbf{可选基础补课 · 2026-09-23}


\StudyRepoRef{CS5222/course-notes/README.md}{}{CS5222目录} · \StudyRef{Foundations-MathForML}{foundation-nav}{数学补课}


网络计算最常见的绊脚石，不是公式复杂，而是把字节当比特、把距离当速率、把“第一个到”当“整个都到了”。这份补课只拆这些台阶，协议机制仍在 Chapter 正文讲。小算例均为教学补充。


\begingroup
\small\setlength{\tabcolsep}{3pt}
\renewcommand{\arraystretch}{1.22}
\begin{longtable}{@{}>{\raggedright\arraybackslash}p{\dimexpr0.324049\linewidth-1.33333\tabcolsep\relax}>{\raggedright\arraybackslash}p{\dimexpr0.262114\linewidth-1.33333\tabcolsep\relax}>{\raggedright\arraybackslash}p{\dimexpr0.413837\linewidth-1.33333\tabcolsep\relax}@{}}
\toprule
卡在哪里 & 入口 & 自查（答案在对应小节） \\
\midrule\endfirsthead
卡在哪里 & 入口 & 自查（答案在对应小节） \\
\midrule\endhead
\bottomrule\endfoot
bit、byte、Mbps分不清 & \StudyRef{Foundations-NetworkBasics}{units}{单位与数量级} & 1500 B以2 Mbps发出要多久？ \\
二进制和IP字段陌生 & \StudyRef{Foundations-NetworkBasics}{binary}{二进制与字节} & 00000101表示多少？ \\
英文报文看起来像乱码 & \StudyRef{Foundations-NetworkBasics}{messages}{报文、字段与分隔} & CRLF空行起什么作用？ \\
时间轴不知道从哪里数 & \StudyRef{Foundations-NetworkBasics}{timeline}{第一位、最后一位和流水线} & 2 ms发送、8 ms传播，最后一位何时到？ \\
\end{longtable}
\endgroup






\StudyHeading{1}{1. Bit, byte and rate｜仓库有多少货，和每秒搬多少货，是两件事}
\StudyTarget{Foundations-NetworkBasics__units}
\StudyTarget{Foundations-NetworkBasics__1-bit-byte-and-rate}


bit（比特）是一个二进制位，可取0或1。byte（字节）通常是8 bit，本课程采用8位字节。大写B与小写b不能省略：1500 B=12000 bit；文件大小是数量，2 Mbps是每秒2000000 bit的速率。


本课程网络速率用十进制：k=10\textsuperscript{3}，M=10\textsuperscript{6}，G=10\textsuperscript{9}。时间中milli是千分之一（ms=10\textsuperscript{-}\textsuperscript{3}s），micro是百万分之一（\ensuremath{\mu}s=10\textsuperscript{-}\textsuperscript{6}s）；Mbps中的大写M和ms中的小写m方向相反。KiB、MiB才明确表示1024、1024\textsuperscript{2}字节；题目中的Mbps按十进制换算。
\[
\text{发送时间}=\frac{\text{数据量(bit)}}{\text{速率(bit/s)}}.\tag{N.1}
\]


\begin{StudyBlock}{完整示范 / Worked example}{先统一bit与秒}
1500 B、R=2 Mbps。先换 \(L=1500\times8=12000\) bit，再算 \(L/R=12000/(2\times10^6)=0.006\) s=6 ms。单位“bit\ensuremath{\div}bit/s”最后剩s，是一个自带报警器：如果还剩米或字节，通常哪里没对齐。
\end{StudyBlock}


传播是另一件事。距离d=1000 km=10\textsuperscript{6}m，速度s=2\ensuremath{\times}10\textsuperscript{8}m/s，\(d/s=0.005\)s=5ms。提高链路发送速率R不等于让电磁信号速度s跟着翻倍。


\StudyImage{learning/foundation-notes/assets/network-units.png}
\StudyCaption{发送与传播使用不同的“数量除以速率”：上条处理比特，下条处理距离。}


上条是“把多少bit送上路”，下条是“已经上路的bit走多远”。两个除法输出同为秒，但分子和分母的含义不同。带宽还可能在物理通信里指Hz频宽；本章说link bandwidth时通常指bit/s容量，要看上下文。


\textbf{变式 / Transfer:} 1000 B以10 Mbps发送，链路2000 km、传播速度2\ensuremath{\times}10\textsuperscript{8}m/s，忽略其他时延，整包到达要多久？ / Send 1000 B at 10 Mbps over 2000 km at propagation speed 2\ensuremath{\times}10\textsuperscript{8} m/s. Ignore other delays and find complete-packet arrival time.



\begin{StudySolution}{答案 / Answer}


发送0.8ms，传播10ms，总10.8ms。 / Transmission is 0.8 ms, propagation is 10 ms, total 10.8 ms. 计算前先把B换成bit、Mbps换成bit/s。


\end{StudySolution}


\textbf{English:} A byte contains eight bits. Convert all quantities before dividing; link rate and propagation speed describe different processes.


返回：\StudyRef{CS5222-Chapter01}{delay}{Chapter1：时延} · \StudyRef{CS5222-Tutorial02}{root}{Tutorial2} · \StudyRef{CS5222-Assignment01}{p6}{Assignment1 P6}






\StudyHeading{1}{2. Binary and octets｜位的位置，也在贡献数值}
\StudyTarget{Foundations-NetworkBasics__binary}
\StudyTarget{Foundations-NetworkBasics__2-binary-and-octets}


十进制305表示3\ensuremath{\times}100+0\ensuremath{\times}10+5\ensuremath{\times}1。二进制每向左一位权重翻倍：一个8位字节的权重是128、64、32、16、8、4、2、1。


\begingroup
\small\setlength{\tabcolsep}{3pt}
\renewcommand{\arraystretch}{1.22}
\begin{longtable}{@{}>{\raggedright\arraybackslash}p{\dimexpr0.111111\linewidth-1.77778\tabcolsep\relax}>{\raggedright\arraybackslash}p{\dimexpr0.111111\linewidth-1.77778\tabcolsep\relax}>{\raggedright\arraybackslash}p{\dimexpr0.111111\linewidth-1.77778\tabcolsep\relax}>{\raggedright\arraybackslash}p{\dimexpr0.111111\linewidth-1.77778\tabcolsep\relax}>{\raggedright\arraybackslash}p{\dimexpr0.111111\linewidth-1.77778\tabcolsep\relax}>{\raggedright\arraybackslash}p{\dimexpr0.111111\linewidth-1.77778\tabcolsep\relax}>{\raggedright\arraybackslash}p{\dimexpr0.111111\linewidth-1.77778\tabcolsep\relax}>{\raggedright\arraybackslash}p{\dimexpr0.111111\linewidth-1.77778\tabcolsep\relax}>{\raggedright\arraybackslash}p{\dimexpr0.111111\linewidth-1.77778\tabcolsep\relax}@{}}
\toprule
位权 & 128 & 64 & 32 & 16 & 8 & 4 & 2 & 1 \\
\midrule\endfirsthead
位权 & 128 & 64 & 32 & 16 & 8 & 4 & 2 & 1 \\
\midrule\endhead
\bottomrule\endfoot
00000101 & 0 & 0 & 0 & 0 & 0 & 1 & 0 & 1 \\
\end{longtable}
\endgroup


只把为1的列相加，4+1=5。8位无符号数最小0、最大255，共256种。IPv4地址用四个这样的八位组（octet）展示，例如192.0.2.1：192=128+64，写为11000000；整个地址是32位。这个例子使用文档示例地址，不是让你访问的真实服务器。


端口号是另一个字段，帮助定位主机上的服务/进程；并不是IP的“第五个点分数字”。域名是人方便使用的名称，经DNS解析得到记录；名称、IP和端口分别回答不同问题。协议更详细的命名、寻址与复用在Chapter2讲。


十六进制（hexadecimal）每位代表4个二进制位，数字0–9后用A–F表示10–15。0x0D=13，0x0A=10，所以能在抓包中用紧凑的两位hex表示一个字节。\StudyInline{0x}是记号，不是额外数据。


\textbf{变式 / Transfer:} 00001010的十进制值是多少？11000000呢？ / Convert these two binary octets to decimal.



\begin{StudySolution}{答案 / Answer}


8+2=10；128+64=192。 / They represent 10 and 192. 前导0不改变数值，却帮助我们看出字段固定8位。


\end{StudySolution}


\textbf{English:} Binary digits are weighted by powers of two. An IPv4 address contains four octets; a port identifies a transport endpoint within a host context.


返回：\StudyRef{CS5222-Chapter01}{core}{Chapter1：转发表示} · \StudyRef{CS5222-Chapter02}{processes}{Chapter2：进程与socket}






\StudyHeading{1}{3. Messages and bytes｜一段报文不是一整块没有边界的文字}
\StudyTarget{Foundations-NetworkBasics__messages}
\StudyTarget{Foundations-NetworkBasics__3-messages-and-bytes}


先看教学HTTP例子，保留英文协议文本：



\begin{StudyCode}
GET /notes.html HTTP/1.1\r\n
Host: example.com\r\n
Connection: close\r\n
\r\n
\end{StudyCode}


第一行是方法、路径、版本；后面是键值头字段。这里展示的\StudyInline{\textbackslash{}r\textbackslash{}n}是两个控制字节的可读写法：CR（carriage return，0x0D）和LF（line feed，0x0A），不是反斜杠、r、反斜杠、n四个普通字节。连续CRLF形成空行，结束头部；不是每个报文都有非空body。


文本字符串与字节序列要区分：Python的\StudyInline{'A'}是文本，\StudyInline{b'A'}是ASCII字节；\StudyInline{'A'.\allowbreak{}encode(\allowbreak{}'utf-8')\allowbreak{}}得到一个字节，而汉字在UTF-8下通常需要多个字节。因此\StudyInline{len(\allowbreak{}text)\allowbreak{}}不总等于网络发送的字节数。原Tutorial3使用ASCII报文，ASCII范围一个字符对应一个字节。


TCP提供有序字节流，应用必须自己定义消息边界；一次\StudyInline{send}不保证对应对方一次\StudyInline{recv}。UDP保留数据报边界，但接收缓冲不够会截断；可靠性与顺序不是默认保证。这里先知道“读到了几个字节”和“收齐一个完整应用消息”可能不同。


\textbf{变式 / Transfer:} \StudyInline{b'OK\textbackslash{}r\textbackslash{}n'}有多少字节？报文头中的User-Agent能否证明实际浏览器品牌？ / How many bytes are in b'OK\textbackslash{}r\textbackslash{}n', and does User-Agent prove the actual browser identity?



\begin{StudySolution}{答案 / Answer}


4字节：O、K、CR、LF。User-Agent只是客户端提供的字段，可以修改；只能说明报文自称是什么。 / Four bytes; User-Agent is a client-supplied claim, not proof of the actual browser.


\end{StudySolution}


\textbf{English:} Preserve protocol delimiters exactly. Application messages need framing, especially over a TCP byte stream.


返回：\StudyRef{CS5222-Chapter02}{http}{Chapter2：HTTP报文} · \StudyRef{CS5222-Chapter03}{tcp}{Chapter3：Telnet字节序号} · \StudyRef{CS5222-Tutorial03}{q1}{Tutorial3 Q1} · \StudyRef{CS5222-Assignment01}{pop3}{Assignment1 POP3}






\StudyHeading{1}{4. Timelines and pipelines｜先问“谁到哪了”，再问“用了多久”}
\StudyTarget{Foundations-NetworkBasics__timeline}
\StudyTarget{Foundations-NetworkBasics__4-timelines-and-pipelines}


规定t=0为发送第一位的起点。一整包需要T=2ms送入链路，单个位在链路中的传播时间为\ensuremath{\tau}=8ms。连续流体近似下：第一位t=0开始走，t=8ms到；最后一位t=2ms离开发送端，t=10ms到。整包到齐需T+\ensuremath{\tau}，而不是只看首位到达。


\StudyImage{learning/foundation-notes/assets/network-timeline.png}


\StudyCaption{沿图横轴读时间，纵轴区分发送端与接收端。两条斜线是首/末位轨迹，间距是发送时长；线上的位置表示比特在该时刻到达哪里。如果\ensuremath{\tau}<T，首位可以在末位离开发送端之前到达，对网络并无矛盾。}


存储转发（store-and-forward）要求中间节点收到整包再转下一条链路。若忽略传播/排队/处理，两条相同速率链路、3个相同长度包，每包发送时间T：


\begingroup
\small\setlength{\tabcolsep}{3pt}
\renewcommand{\arraystretch}{1.22}
\begin{longtable}{@{}>{\raggedright\arraybackslash}p{\dimexpr0.333333\linewidth-1.33333\tabcolsep\relax}>{\raggedright\arraybackslash}p{\dimexpr0.333333\linewidth-1.33333\tabcolsep\relax}>{\raggedright\arraybackslash}p{\dimexpr0.333333\linewidth-1.33333\tabcolsep\relax}@{}}
\toprule
时间段 & 第1条链路 & 第2条链路 \\
\midrule\endfirsthead
时间段 & 第1条链路 & 第2条链路 \\
\midrule\endhead
\bottomrule\endfoot
0–T & 包1 & 等待 \\
T–2T & 包2 & 包1 \\
2T–3T & 包3 & 包2 \\
3T–4T & 空闲 & 包3 \\
\end{longtable}
\endgroup


总4T，不是每包独自跑完两跳再发下一个所得6T，因为不同链路能同时处理不同包。等长包、等速率、无其他等待时，P个包经过N条链路是 \((P+N-1)T\)；有不同速率、头部开销或其他流量时要重新画时间轴。


\textbf{变式 / Transfer:} 同样条件，4个包、3条链路，每包发送2ms，总多久？ / Four packets cross three equal-rate links, with 2 ms transmission per packet and no other delays. Find the completion time.



\begin{StudySolution}{答案 / Answer}


\((4+3-1)\times2=12\)ms。 / 12 ms. 第一包用3T，第2–4包每隔T到达，合计6T。


\end{StudySolution}


\textbf{English:} Name the event before calculating its time. Pipelining overlaps work on different links; store-and-forward still requires each individual packet to arrive completely.


返回：\StudyRef{CS5222-Chapter01}{switching}{Chapter1：存储转发} · \StudyRef{CS5222-Chapter01}{queue-events}{Chapter1：队列事件} · \StudyRef{CS5222-Chapter03}{utilization}{Chapter3：停等利用率} · \StudyRef{CS5222-Tutorial01}{q1}{Tutorial1 Q1} · \StudyRef{CS5222-Assignment01}{p31}{Assignment1 P31}



\StudyHeading{1}{来源与复算}
\StudyTarget{Foundations-NetworkBasics___1}


这些先修概念直接服务当前Chapter1 pp25–26、45–48、63–65，Chapter2进程、HTTP与socket，以及当前Tutorial/Assignment。可按对应课程中的需要选读。示意图与数值由 \StudyRepoRef{CS5222/course-notes/tools/network_examples.py}{}{network\_examples.py} 生成，结果在\StudyRepoRef{CS5222/course-notes/network-example-results.json}{}{复算记录}。组合概率另见\StudyRef{Foundations-MathForML}{combinations}{MathForML §6}。
